\documentclass[a4paper,11pt]{jsarticle}
\usepackage{amsmath}
\usepackage{amssymb}
\usepackage{bm} % ベクトル太字用

\begin{document}

\section{異方性を考慮した視線積分モデル}

本研究では、計測器で観測される輝度 $S$ を、視線（Line of Sight: LOS）に沿った放射強度 $\varepsilon$ の積分として計算する。ここで、従来の等方的な放射モデルとは異なり、局所的な磁場ベクトル $\vec{B}$ と視線ベクトル $\vec{l}$ のなす角 $\theta$ に依存する放射の異方性を導入する。

観測される輝度 $S$ は、以下の線積分によって定式化される。

\begin{equation}
    S = \int_{L} \varepsilon(R, Z) \cdot \mathcal{F}(\theta) \, dl
\end{equation}

ここで、$L$ は視線の積分経路、$dl$ は視線上の微小距離要素である。$\varepsilon(R, Z)$ は磁気面形状等に依存する空間的な放射強度分布（ファントム）を表し、$\mathcal{F}(\theta)$ は磁場方向に対する放射の異方性係数（正規化された放射パターン）を表す。

\subsection{座標変換と磁場-視線角 $\theta$ の計算}

視線上の任意の位置 $\vec{r} = (x, y, z)$ における磁場ベクトル $\vec{B}$ の向きを正確に評価するため、円筒座標系 $(R, \phi, Z)$ で与えられた平衡磁場データ $(B_R, B_T, B_Z)$ を、視線ベクトルが定義されている実験室直交座標系 $(x, y, z)$ へと変換する。

視線上の点におけるトロイダル角を $\phi = \arctan(y/x)$ とすると、直交座標系での磁場ベクトル $\vec{B}_{\mathrm{xyz}}$ は、以下のような回転行列を用いて記述される。

\begin{equation}
    \vec{B}_{\mathrm{xyz}} = 
    \begin{pmatrix}
        B_x \\
        B_y \\
        B_z
    \end{pmatrix}
    =
    \begin{pmatrix}
        \cos\phi & -\sin\phi & 0 \\
        \sin\phi & \cos\phi & 0 \\
        0 & 0 & 1
    \end{pmatrix}
    \begin{pmatrix}
        B_R \\
        B_T \\
        B_Z
    \end{pmatrix}
\end{equation}

これより、$x$成分および$y$成分はそれぞれ $B_x = B_R \cos\phi - B_T \sin\phi$、$B_y = B_R \sin\phi + B_T \cos\phi$ となる。視線の単位ベクトルを $\vec{n}_{\mathrm{los}}$ とすると、局所磁場ベクトルとのなす角 $\theta$ は内積を用いて次式で算出される。

\begin{equation}
    \cos\theta = \frac{\vec{B}_{\mathrm{xyz}} \cdot \vec{n}_{\mathrm{los}}}{|\vec{B}_{\mathrm{xyz}}|}
\end{equation}

シミュレーションでは、この $\theta$ を用いて、事前に計算されたルックアップテーブル（LUT）から異方性係数 $\mathcal{F}(\theta)$ を参照する。

\subsection{微視的放射モデル}

異方性係数 $\mathcal{F}(\theta)$ は、特定のピッチ角 $\alpha_e$ を持つ電子群からの制動放射（Bremsstrahlung）をモデル化して導出される。電子の速度ベクトル $\vec{v}$ が磁力線に対してピッチ角 $\alpha_e$ でジャイロ回転していると仮定し、観測方向に対する放射強度をジャイロ位相 $\varphi_{\mathrm{gyro}}$ について平均化する。

単一電子からの放射強度分布 $I(\psi)$ が、双極子放射と相対論的ビーミング効果によって記述されると仮定すると、角度 $\theta$ 方向への放射強度 $\mathcal{F}(\theta)$ は以下のように積分形式で表される。

\begin{equation}
    \mathcal{F}(\theta) \propto \int_{h\nu_{\min}}^{h\nu_{\max}} T(h\nu) \left[ \frac{1}{2\pi} \int_{0}^{2\pi} I(\psi(\theta, \alpha_e, \varphi_{\mathrm{gyro}})) \, d\varphi_{\mathrm{gyro}} \right] d(h\nu)
\end{equation}

ここで、$T(h\nu)$ はフィルタの透過率関数、$\psi$ は電子の瞬間的な速度ベクトルと視線のなす角である。本モデルにより、磁力線方向に対する電子の運動の異方性が、観測される放射強度の角度依存性に反映される。

\end{document}